\section{Intra-node performance}
\label{section:tailoring:intra-node}
As referenced inside Section~\ref{section:high-level:intra-node}, the program is considered to be making good use of a system if it is able to use more than 80\% of its cores with more than 80\% of (strong scaling) efficiency. This metric is not workload agnostic. For instance, a program could have two benchmarks with vastly different problem sizes. Additionally, for programs with simple scaling in problem size, for example scaling a system size up resulting in a quadratic increase in compute time, the ratio of memory usage to compute cannot be assumed to relate similarly. 

In order to evaluate the intra-node performance of a code therefore, choices must be made as to a reasonable range of workloads for the program to perform. The range of workloads should aim to be representative of real world behaviour of the program, and as such will require information about the program's use cases to be effective. For intra-node performance, this will often relate to program's utilising workstations, or batch parallel problems on beowulf clusters. For problems which will exclusively be run on multiple nodes due to their size or computation length, intra-node scaling can be difficult to treat. In these cases it may be necessary to examine a small number of cores on a fixed number of nodes and gradually increase the number of cores. This may be prohibitively time consuming however, and should generally only be examined where necessary for larger problems. 

The range of workloads for different codes could include:
\begin{itemize}
		\item A large workload, with a large memory footprint
		\item A compute heavy workload, with fixed memory
	        \item A workload with lots of memory accesses
		\item An I/O heavy workload 
\end{itemize}

Unlike inter-node scaling intra-node scaling generally won't exhibit super scaling behaviour (where the performance increases more than proportionally with core count). In inter-node scaling that is usually because of additional resources available across nodes, for example increasing the cache space, or i/o resources. 

\subsection{Configuration}
Where possible, the configuration of a node should be comparable to the expected hardware configurations in wider use to run the program. This may point to a reduction in the number of cores, or available memory for a node. It can be important, however, to measure the performance to a full node if possible, in order to evaluate the benefits of larger hardware for future procurement.

For heterogenous computing, there is at least a three dimensional state space to examine. The number of cores per node, the number of GPU's per node, and the class of GPU can all be important variables, and the total cost for procuring and running some combination of these can be an important metric. For instance for a machine with nodes configured with different numbers of different GPU's an example configuration table could look like:

% Please add the following required packages to your document preamble:
% \usepackage{multirow}
\begin{table}[]
\begin{tabular}{llllllll}
\multirow{2}{*}{ID} & \multirow{2}{*}{Cores} & \multirow{2}{*}{GPUs} & \multirow{2}{*}{GPU Class} & \multicolumn{2}{l}{Performance Metric 1} & \multicolumn{2}{l}{Performance Metric 2} \\
                    &                        &                       &                            & Benchmark 1         & Benchmark 2        & Benchmark 1         & Benchmark 2        \\
                    &                        &                       &                            &                     &                    &                     &                    \\
                    &                        &                       &                            &                     &                    &                     &                   
\end{tabular}
\end{table}


\subsection{Metrics}
Metrics for measuring intra-node performance

\subsection{Pinning}
\label{section:tailoring:intra-node:pinning}

% I think we introduce here the likwid-pin and likwid-mpirun interfaces properly. We demonstrate a workflow of: (1) look at likwid-topology and think of different ways to configure your ranks, e.g., by node? by socket? by NUMA domain? (2) Then if there are threads, how do we include them: spread or close? We can be highly explicit in this with LIKWID. (3) We run experiments and test out different ways of distributing. We might want to try testing things all packed together vs. spread out.

% always check debugging information to see the placement of the threads and ranks

One external dial that we can tweak to understand a code's intra-node
performance is the affinity and placement of processes. Whether this be MPI
ranks or threads from a multithreading library, such as OpenMP, we have
flexibility in where these resources are placed. Some codes come with
recommendations such as `assign 1 MPI rank per socket` but the architectures of
individual chips is always evolving hence it worth revisiting this.

We cannot be entirely exhaustive in experimenting with thread placement
patterns, and neither is this worth while. For hybrid parallel applications,
e.g., MPI+OpenMP, we can have a large parameter space of varying rank count vs
thread count. To reduce this search we instead recommend a simplified workflow
to just begin testing resource allocation for intra-node performance.

\begin{checklist}
 \item Remind yourself, from the pre-assessment, of the structure of the
	 hardware that you are running on. In particular, the domains of the
	 node: size of socket, NUMA domains, and even chiplets or other
	 substructures.
 \item Compare with your high-level assessment and plan out which
	 configurations have not been explored.
 \item Run strong scaling analyses of these different configurations, and
	 calculate the high-level intra-node metric from
	 Section~\ref{section:high-level:intra-node}. 
 \item Note down any findings, such as better or worse performance in different
	 configurations, in an assessment report.
\end{checklist}

This checklist might initially seem vague but we can walk through a simple
example. Say we are tasked with analysing a MPI+OpenMP code for which we have
been advised to use one 1 MPI rank per socket, but the submitter has specified
they are running on a chip which happens to have many small NUMA domains. We
would run through our high-level assessment using what the submitter
recommends, but during this tailoring stage we believe it is worth testing two
different cases: 1 rank per node, and 1 rank per NUMA domain. 


\begin{instructions}{\likwid}
	TBD
\end{instructions}

What about thread spread or close?



